\documentclass[12pt,a4paper]{article}
\usepackage[utf8]{inputenc}
\usepackage[margin=1in]{geometry}
\usepackage{amsmath}
\usepackage{amssymb}
\usepackage{graphicx}
\usepackage{xcolor}
\usepackage{hyperref}
\usepackage{fancyhdr}
\usepackage{array}
\usepackage{booktabs}

\pagestyle{fancy}
\lhead{Estratégia Adaptativa Multi-Fator}
\rhead{Etapa 1: Tese}
\cfoot{\thepage}

\title{\textbf{TESE ESTRATÉGICA}\\[0.5cm]
\textit{Adaptive Multi-Factor Regime-Switching}}
\author{Quantitative Research}
\date{\today}

\begin{document}

\maketitle

\section*{Resumo Executivo}

A presente estratégia propõe uma abordagem sofisticada de alocação dinâmica de capital através da combinação de três pilares: (i) detecção de regime de mercado via Hidden Markov Model, (ii) exposição a múltiplos fatores quantificados, e (iii) ensemble machine learning para previsão de retornos. A hipótese central é que mercados não são homogêneos—diferentes regimes (bull, bear, high-volatility) requerem alocações distintas para otimizar o tradeoff risco-retorno.

\begin{figure}[h]
\centering
\includegraphics[width=0.95\textwidth]{reports/build/model_pipeline.png}
\caption{Estrutura operacional da estratégia adaptativa multi-fator}
\end{figure}

\section{Fundamento Teórico}

\subsection{Hipótese 1: Multi-Fator de Fama-French}

Baseamo-nos na extensão empírica do CAPM proposta por Fama \& French (2015). Retornos esperados são explicados por exposição a fatores sistêmicos:

\[
E[R_i] = R_f + \beta_{market} \cdot \lambda_{market} + \beta_{value} \cdot \lambda_{value} + \beta_{momentum} \cdot \lambda_{momentum} + \varepsilon_i
\]

Onde:
\begin{itemize}
    \item $\beta_{value}$: exposição ao fator VALUE (high book-to-market)
    \item $\beta_{momentum}$: exposição ao fator MOMENTUM (retornos passados de 12m)
    \item $\lambda$: prêmio de risco associado a cada fator
    \item $\varepsilon_i$: alfa específico do ativo
\end{itemize}

\textbf{Intuição}: Ativos com elevado book-to-market oferecem maior retorno esperado (fator value). Ativos com momentum positivo tendem a continuar subindo (fator momentum).

\subsection{Hipótese 2: Detecção de Regime via HMM}

Markets exibem diferentes regimes de volatilidade e tendência. Modelamos isso como um processo Markoviano latente com 3 estados:

\[
s_t \in \{\text{Bull}, \text{Bear}, \text{HighVol}\}
\]

Cada estado é caracterizado por distribuição específica de retornos e volatilidade:

\[
r_t \mid s_t = i \sim \mathcal{N}(\mu_i, \sigma_i^2)
\]

\begin{table}[h]
\centering
\begin{tabular}{lccc}
\toprule
\textbf{Regime} & $\mu_i$ (retorno médio) & $\sigma_i$ (volatilidade) & Duração típica \\
\midrule
Bull & $+0.08\% / \text{dia}$ & $0.8\%$ & 3-6 meses \\
Bear & $-0.06\% / \text{dia}$ & $1.2\%$ & 2-4 meses \\
HighVol & $+0.02\% / \text{dia}$ & $2.0\%$ & 1-2 meses \\
\bottomrule
\end{tabular}
\end{table}

\subsection{Hipótese 3: Machine Learning para Previsão}

Combina dois modelos para captar relações não-lineares entre fatores e retornos:

\subsubsection{Gradient Boosting}
Modelo não-paramétrico que iterativamente reduz erros residuais. Captura interações complexas entre fatores.

\subsubsection{Ridge Regression}
Modelo linear regularizado. Oferece estabilidade e interpretabilidade. Complementa GBM em regimes com relações lineares.

\textbf{Ensemble Weights}:
\[
\hat{r}_t = w_{GB} \cdot \hat{r}_{GB,t} + w_{Ridge} \cdot \hat{r}_{Ridge,t}
\]

Onde pesos $(w_{GB}, w_{Ridge})$ variam conforme regime detectado para maximizar sharpe ratio esperado.

\subsection{Hipótese 4: Otimização Dinâmica}

Portfolio é otimizado usando mean-variance com restrições dinâmicas:

\[
\max_w \quad w^T \hat{\mu} - \frac{\lambda}{2} w^T \Sigma w
\]

Sujeito a:
\begin{align}
\sum_i w_i &\leq 1 + \text{leverage}(s_t) \\
|w_i| &\leq w_{max}(s_t) \\
|\Delta w| &\leq \text{turnover}_{\max}
\end{align}

\textbf{Leverage dinâmico por regime}:
\begin{itemize}
    \item Bull: leverage = 1.3 (tomar risco adicional)
    \item Bear: leverage = 0.7 (reduzir exposição)
    \item HighVol: leverage = 1.0 (posição neutra)
\end{itemize}

\section{Justificativa Teórica}

\subsection{Por que múltiplos fatores?}

Fatores individuais são instáveis. Combinação reduz risco idiossincrático:
\[
\text{Volatilidade idiossincrática} \propto \frac{1}{\sqrt{n_{\text{fatores}}}}
\]

\subsection{Por que regime detection?}

Regimes definem estrutura de covariância:
\[
\Sigma_t = f(s_t) \implies \text{Otimização depende de } s_t
\]

Durante market crashes, correlações explodem. HMM antecipa regime antes de ser tarde demais.

\subsection{Por que ensemble ML?}

Relação entre fatores e retornos muda conforme regime. Ensemble capta:
\begin{itemize}
    \item Linearidade em regimes estáveis (Ridge)
    \item Não-linearidade em inflexões de regime (GBM)
\end{itemize}

\section{Vantagens Competitivas}

\begin{enumerate}
    \item \textbf{Adaptativo}: Ajusta agressividade conforme regime, evitando over-leverage em crashes
    \item \textbf{Diversificado}: 7 fatores independentes reduzem tail-risk
    \item \textbf{Interpretável}: HMM + fatores definem decisões explícitas (não black-box)
    \item \textbf{Robusto}: Ensemble + regularização reduzem overfitting
    \item \textbf{Escalável}: Aplicável a qualquer universo de ativos (ações, futuros, cripto)
\end{enumerate}

\section{Risks e Limitações}

\begin{enumerate}
    \item \textbf{Parameter Estimation Risk}: Estimação de $\Sigma$ requer suficiente histórico
    \item \textbf{Regime Misclassification}: HMM pode ser enganado por outliers
    \item \textbf{Look-Ahead Bias}: Uso de fatores fundamentalistas pode sofrer delay de publicação
    \item \textbf{Overfitting em ML}: GBM requer validação cruzada rigorosa
    \item \textbf{Model Risk}: Pressupostos teóricos podem não valer em stress scenarios
\end{enumerate}

\section{Próximas Etapas}

A presente tese será validada através de:
\begin{enumerate}
    \item Implementação do modelo em Python
    \item Backtest em 3 anos de dados históricos
    \item Walk-forward analysis e stress testing
    \item Comparação com benchmarks (S\&P 500, 60/40 portfolio)
\end{enumerate}

\end{document}
